%=========================================================
% APPENDICE
%=========================================================
\appendix
\chapter[APPENDICE]{Codice sorgente}

Questa appendice raccoglie il codice sorgente principale relativo alle componenti
implementative del protocollo di autenticazione basato su Schnorr.  
Le classi riportate rappresentano i moduli centrali per la gestione delle chiavi crittografiche
e per l'esecuzione delle fasi di prova e verifica all'interno del protocollo.

%---------------------------------------------------------
\section{\texttt{KeyManager}}
\label{sec:key_manager}
%---------------------------------------------------------

La classe \texttt{KeyManager} gestisce la memorizzazione sicura delle chiavi private
degli utenti. Se viene fornita una password, la chiave viene cifrata con AES-GCM
utilizzando una chiave derivata tramite Scrypt. In assenza di password, la chiave
viene salvata in chiaro nel file system locale.

\begin{listing}
\begin{minted}[
    frame=lines,
    linenos,
    breaklines,
    fontsize=\fontsize{9pt}{11pt}\selectfont
]{python}
import os
import cryptography
from pathlib import Path
from cryptography.hazmat.primitives.kdf.scrypt import Scrypt
from cryptography.hazmat.primitives.ciphers.aead import AESGCM

class KeyManager:
    HOME_PATH = Path.home()
    SCHNORR_DIR = HOME_PATH / ".schnorr"

    @staticmethod
    def _derive_key(password: str, salt: bytes) -> bytes:
        """Deriva una chiave AES-256 (32 byte) usando PBKDF2-HMAC-SHA256"""
        ...
    
    @classmethod
    def save_private_key(cls, username: str, key: int, password: str = None) -> None:
        cls.SCHNORR_DIR.mkdir(parents=True, exist_ok=True)
        privkey_path = cls.SCHNORR_DIR / f"{username}_privkey.pem"

        if not password:
            cls._save_private_key(username, key)
            return

        data = str(key).encode()
        salt, nonce = os.urandom(16), os.urandom(12)
        aesgcm = AESGCM(cls._derive_key(password, salt))
        ciphertext = aesgcm.encrypt(nonce, data, None)
        with open(privkey_path, "wb") as f:
            f.write(salt + nonce + ciphertext)

    @classmethod
    def load_private_key(cls, username: str, password: str = None) -> int:
        privkey_path = cls.SCHNORR_DIR / f"{username}_privkey.pem"

        if not password:
            return cls._load_private_key(username)

        with open(privkey_path, "rb") as f:
            blob = f.read()
        salt, nonce, ciphertext = blob[:16], blob[16:28], blob[28:]
        aesgcm = AESGCM(cls._derive_key(password, salt))
        plaintext = aesgcm.decrypt(nonce, ciphertext, None)
        return int(plaintext.decode())
\end{minted}
\caption{Implementazione della classe \texttt{KeyManager}.}
\end{listing}

%---------------------------------------------------------
\section{\texttt{SchnorrProver} e \texttt{SchnorrVerifier}}
\label{sec:schnorr_classi}
%---------------------------------------------------------

Le classi \texttt{SchnorrProver} e \texttt{SchnorrVerifier} implementano rispettivamente
il ruolo del Prover e del Verifier nel protocollo di identificazione di Schnorr.
La prima gestisce la generazione delle chiavi e la produzione della risposta,
mentre la seconda verifica la validità della prova o della firma associata.

\begin{listing}
\begin{minted}[
    frame=lines,
    linenos,
    breaklines,
    fontsize=\fontsize{8.5pt}{10pt}\selectfont
]{python}
class Schnorr:
    def __init__(self, group_id: GroupType):
        self.crypto_group = Rfc3526(group_id)
        self.p, self.g, self.q = self.crypto_group.get_params()
        self.public_key, self.public_key_temp = None, None

    def compute_challenge(self, message: int) -> int:
        data = str(self.public_key_temp) + str(self.public_key) + message
        return int(hashlib.sha256(data.encode()).hexdigest(), 16) % self.q

class SchnorrProver(Schnorr):
    def __init__(self, group_id: GroupType):
        super().__init__(group_id, filename)
        self.alpha, self.alpha_temp = None, None
        
    def public_key(self) -> int:
        if self.public_key is None:
            self.public_key = pow(self.g, self.alpha, self.p)
        return self.public_key

    def public_key_temp(self) -> int:
        self.alpha_temp = random.randint(1, self.q - 1)
        self.public_key_temp = pow(self.g, self.alpha_temp, self.p)
        return self.public_key_temp

    def gen_keys(self):
        self.alpha = random.randint(1, self.q - 1)
        self.public_key = pow(self.g, self.alpha, self.p)

    def response(self, challenge: int) -> int:
        return (self.alpha_temp + self.alpha * challenge) % self.q

    def sign_message(self, message: str) -> dict[str, int]:
        challenge = self.compute_challenge(message)
        sign = {"public_key_temp": self.public_key_temp, "response": self.response(challenge)}
        return sign
        
class SchnorrVerifier(Schnorr):
    def __init__(self, group_id: GroupType):
        super().__init__(group_id, filename)
        self.challenge = random.randint(0, self._q - 1)

    def check(self, response: int) -> bool:
        left = pow(self.g, response, self.p)
        right = (self.public_key_temp * pow(self.public_key, self.challenge, self.p)) % self.p
        return left == right
    
    def verify_sign(self, sign: dict, message: str) -> bool:
        self.public_key_temp = sign["public_key_temp"]
        self.challenge = self.compute_challenge(message)
        return self.check(sign["response"])
\end{minted}
\caption{Classi \texttt{SchnorrProver} e \texttt{SchnorrVerifier}.}
\end{listing}

\clearpage
